% !TEX root = ../../main/main.tex
\hypearticle
{O Papel dos Modelos de Mundo na Próxima Era da IA Física}
{Marcelo Trevine}
{

\subtitle{A Imaginação como Superpoder}

\hspace*{1.5em}A mente humana possui a incrível capacidade de prever o que ainda não aconteceu. Quando olhamos para a foto de um saltador olímpico suspenso no ar, projetamos instantaneamente o que está fora do enquadramento visual: inferimos a presença de uma piscina cheia de água e antecipamos o mergulho iminente. Se a água não estivesse lá, reconheceríamos o risco de um desfecho trágico antes mesmo de o atleta tocar o chão.

\hspace*{1.5em}Esse tipo de raciocínio — que transita da observação passiva para a simulação mental de hipóteses contrafactuais (o famoso "e se...?") — é o verdadeiro motor da inteligência. Atualmente, a Inteligência Artificial passa por uma revolução idêntica: ela está deixando de apenas reagir ao presente para construir um "simulador interno" capaz de antecipar o próximo acontecimento. Mas, para prever o mundo, a IA precisa primeiro aprender a representá-lo. É aqui que o campo se divide em dois caminhos.

\subtitle{Simular vs. Representar}

\hspace*{1.5em}Para estruturar esse simulador interno, a engenharia de IA consolidou duas abordagens principais. De um lado, os modelos de mundo explícitos funcionam como motores de simulação ativa. Eles projetam o futuro passo a passo no tempo para guiar o planejamento e as decisões. Arquiteturas como o Dreamer e o MuZero operam nessa lógica, calculando incertezas e visualizando trajetórias dentro de seu espaço latente — o "espaço de pensamento" matemático da rede neural — antes de executar qualquer movimento.

\hspace*{1.5em}Por outro lado, os modelos de mundo implícitos não precisam gastar tanto processamento para criar ou interpretar imagens. Em vez de tentar prever o futuro por meio de vídeos, eles aprendem conceitos como senso comum e leis da física a partir de representações matemáticas. Modelos como a família JEPA, o Genie e os grandes modelos de linguagem (LLMs) mostram que ideias como gravidade, espaço e continuidade podem ser aprendidas de forma autossupervisionada. No fim, seja criando imagens ou organizando conceitos, uma teoria só faz sentido quando consegue representar o funcionamento do mundo real.

\articleimage{1.0}{../sections/physical-ai/imagens/imagem 1 hype.png}{Robô humanóide e um balão de pensamento com a representação do modelo de mundo explícito à esquerda e implícito à direita.}

\subtitle{IA Física no Mundo Real}

\hspace*{1.5em}Levar a inteligência artificial para o nosso cotidiano exige que ela opere sob as restrições de um ambiente físico repleto de atrito, latência e imprevistos — o chamado sim-to-real gap, a lacuna entre a simulação virtual perfeita e a bagunça do mundo real. Na robótica, essa transição exige um equilíbrio entre dois sistemas: a velocidade do Modo-1 (controle reativo e reflexo direto para loops rápidos de percepção e ação) e a profundidade do Modo-2 (imaginação deliberada, em que o robô prevê as consequências de suas escolhas antes de agir.)

\hspace*{1.5em}Essa busca pelo "cérebro" ideal dos robôs tornou-se o epicentro de uma corrida acelerada na China. Embora a indústria chinesa de humanóides tenha dominado a fabricação de corpos mecânicos complexos, dar às máquinas a capacidade de navegar em cenários imprevisíveis continua sendo o maior gargalo. Em feiras de tecnologia em Xangai, centenas de robôs já preparam café, jogam boxe ou classificam objetos em ambientes controlados, mas falham quando expostos ao caos do mundo aberto.

\hspace*{1.5em}A mesma urgência atinge a condução autônoma, em que os veículos precisam resolver cenários de "cauda longa" (long-tail events) — situações raras e perigosas no trânsito para as quais quase não há dados de treino. Para tomar decisões seguras, veículos podem simular futuros alternativos usando modelos generativos estruturados ou extrair mapas tridimensionais limpos a partir de sensores LiDAR (que prevê regiões ocultas do trânsito sem precisar gerar imagens).

\articleimage{1.0}{../sections/physical-ai/imagens/Imagem 2 Hype.png}{Representação de robô humanóide sendo treinado em ambiente virtual para aprender a manipular objetos no mundo real.}

\hspace*{1.5em}Como a coleta de dados físicos reais é lenta e perigosa, o setor passou a recorrer a plataformas abertas como o NVIDIA Cosmos — um modelo fundacional treinado para prever pixels, física e ações —, integrado a gêmeos digitais no NVIDIA Omniverse, validando robôs e carros em cenários virtuais fotorrealistas antes que eles toquem o asfalto. Apesar desses avanços, a engenharia atual ainda esbarra no mesmo limite que nos separa da Inteligência Geral.

}

\highlight{
    Sobre o NVidia Cosmos:
}

\begin{hypeitemize}
\item \href{https://www.youtube.com/watch?v=9Uch931CdX8}{Vídeo no YouTube};
\end{hypeitemize}

\originalsource{https://arxiv.org/pdf/2606.12783}
